\documentclass[minion]{homework}

% Some common macros
\newcommand{\Reals}{\ensuremath{\mathbb R}}% Gives you a shortcut for writing the blackboard R for the real numbers - \RR
\newcommand{\Rats}{\ensuremath{\mathbb Q}} % Gives you a shortcut for writing the blackboard N for the natural numbers - \NN
\newcommand{\Nats}{\ensuremath{\mathbb N}} % Gives you a shortcut for writing the blackboard N for the natural numbers - \NN
\newcommand{\Ints}{\ensuremath{\mathbb Z}} % Gives you a shortcut for writing the blackboard Z for the integer numbers - \ZZ

\doclabel{Math F404: Homework 11}
\docdate{Due: November 17, 2026}

\begin{document}
\begin{problems}

\problem Suppose $f:X\to Y$ is a bijection. Show that if $A\subset X$, $(f^{-1})^{-1}(A)=A$.
That is, show that the preimage of $A$ under $f^{-1}$ is $A$.  Note that the 
two superscripts $-1$ are doing different work, so they do not simply cancel 
as they would if both meant ``inverse function''.

\problem Let $X$ be a topological space and let $A\subset B\subset X$.
Show that the subspace topology on $A$ as a subspace of $B$ is the same
as the subspace topology on $A$ as a subspace of $X$.  Note that $B$ 
in this problem has the subspace topology (as a subspace of $X$).

\problem Suppose $X$ is a topological space and that $A\subset X$.
Show that $A$ is a compact subset of $X$ if and only if $A$ is a compact 
subset of $A$ (with the subspace topology on $A$).

\problem Let $X$ and $Y$ be topological spaces and let $A\subset X$ and $B\subset Y$.
Then $A\times B$ is a subset of $X\times Y$ and has the subspace topology.
But also, $A$ has the subspace topology (as a subspace of $X$) and similarly 
$B$ has the subspace topology, so $A\times B$ has the product topology 
as a product of two subspaces.  Show that these two topologies are the same.

\problem \textbf{(AI Free)}
Suppose $X$ and $Y$ are topological spaces and $Y$ is compact. Let $x\in Y$
and suppose $W$ is an open set containing the fiber $\{x\}\times Y$.
Show that there is an open set $V\subset X$ with $x\in V$ and $V\times Y\subseteq W$.
Hint: cover the fiber with open rectangles, each contained in $W$.  Show that
you can, in fact, use finitely many of the rectangles.  Then build $V$.

\problem AI Use. Please describe specifically how you used AI to help complete this assignment.

\end{problems}

\end{document}
